\section{Related Work}

{\parskip=2pt
\noindent\textbf{Video Generation for Robotics}: Generating videos conditioned on additional constraints (i.e., language, gesture, trajectory) has made great progress in the robotics community, covering simulation, planning and generalist policy learning~\cite{vlp,thisnthat,robodreamer, unipi, gr2, unisim,genie,gen2act,irasim}. 
Nevertheless, their generated videos are still restricted in the short horizon domain, limiting the capability to cope with complicated instructions. Works such as 
\cite{unisim,genie} predict future robotic actions and states using interactive videos, which serve as the robotic world model for next-generation simulation.
To enhance physical plausibility, studies such as \cite{thisnthat,irasim} incorporate gestures and trajectory control into video generation. 
\cite{vlp,unipi} utilize video diffusion model to search actionable video plans for goal-conditioned policy learning, demonstrating the potential of using video as an intermediate representation for long-horizon planning. Recent works~\cite{gr2,gen2act} resort to generated videos as necessary data blocks for large-scale pretraining, promisingly enabling robots to achieve in-the-wild manipulation tasks. Despite the significant progress made in short-horizon tasks, there is substantial space for improvement in long-term video generation. Considering the challenges in long-horizon video generation, such as spatial and temporal inconsistency, in this work, we explore how to generate a long-horizon multi-task video plan for robotic manipulation tasks with reliable interframe consistency.}

{\parskip=2pt
\noindent\textbf{Video-Driven Policy Learning}: Videos contain rich action information for robotic policy learning. Recent advancements have demonstrated the effectiveness of video-driven policy learning as a promising approach toward developing a generalist policy model~\cite{rdtb,pi_0,diffpolicy,avdc,anypoint}. However, leveraging generated videos to guide robotic learning in long-horizon tasks remains an open challenge. Tracking-based methods~\cite{track2act,avdc,anypoint} rely heavily on keypoint detection and optical flow tracking, making them susceptible to occlusions and degraded video quality. Diffusion Policy~\cite{diffpolicy} learns to act from a denoising process conditioned on video-based observation data, but it still lacks the generalization ability in long-horizon tasks without expert demonstration. VLA models~\cite{rdtb,pi_0} require astronautical-scale pretraining and expensive high-quality data for finetuning.  Similarly, the inverse dynamics model introduced in \cite{unipi} for video-based policy learning suffers from inaccurate estimations when faced with inconsistencies in generated videos. Therefore, to efficiently utilize the generated video, we propose a lightweight but effective transformer-based model to estimate the robot joint state as policy.}
%THIS MIGHT BE ADDED
% Instead of training a policy model to generate an action for each frame, 

{\parskip=2pt
\noindent\textbf{Long-Horizon Video Generation}: 
The length of generated videos is bottlenecked by large computation requirements. The majority of recent works employ autoregressive architecture, inferencing iteratively to compose the long-horizon video from short segments~\cite{opensora,streamingt2v,freenoise,genlvideo,freelong,vdt}. To solve interframe inconsistency, StreamingT2V~\cite{streamingt2v} designs a short memory bank with an appearance preservation module. Freenoise~\cite{freenoise} reschedules the sample noise with motion injection. VDT~\cite{vdt} achieves preserving temporal dependency with spatial-temporal mask
modeling. However, in the robotic manipulation setting with many small objects and large-range motion, it is still challenging to predict long-horizon video plans with good object consistency in number and position. Inspired by recent text-to-video works~\cite{nuwa}, we adopt a hierarchical architecture to reduce noise accumulation and computation expense. However, the nature of our image-to-video task and the more complex data distribution in robotic manipulation tasks impose more challenges of interframe consistency, requiring a more sophisticated architecture design.}
